% 第49--69行建议替换。原有公式标签予以保留。
考虑目标能级 $i$ 与其直接馈入能级 $k$。令 $n_i(t)$ 为时刻 $t$ 的瞬时布居数，$b_{ki}$ 为 $k$ 的全部退激中直接进入 $i$ 的分支比，$\tau_i$ 为平均寿命。在共同的事件归一、近似相同且恒定的飞行速度下，级联率方程为
\begin{equation}\label{eq:5-8-rate-eq}
\frac{\mathrm d n_i}{\mathrm dt}=-\frac{n_i}{\tau_i}+\sum_k b_{ki}\frac{n_k}{\tau_k}.
\end{equation}
定义 $U_i(t)=\int_t^\infty n_i(t')/\tau_i\,\mathrm dt'$ 和 $S_i(t)=\int_0^t n_i(t')/\tau_i\,\mathrm dt'$，它们是该能级全部出口在 $t$ 之后和之前的累计退激数，而非瞬时布居。假设所有相关激发态在积分上限退激，则积分上式得到 $n_i=U_i-\sum_k b_{ki}U_k$；又 $\mathrm dS_i/\mathrm dt=n_i/\tau_i$。在 $t\simeq x/v$ 的低速近似下，
\begin{equation}\label{eq:5-8-ddcm-singles}
\tau_i(x)=\frac{U_i(x)-\sum_k b_{ki}U_k(x)}{v\,\mathrm dS_i(x)/\mathrm dx}.
\end{equation}
$S_i+U_i$ 距离无关时，$\mathrm dS_i/\mathrm dx=-\mathrm dU_i/\mathrm dx$，这一恒等式不要求没有馈入，但寿命式仍须扣除馈入。实际峰面积应按各跃迁的分支、内转换、探测效率、角分布及几何效应转换到相同的母态总退激归一；若直接使用观测曲线，则须保留不同跃迁之间的相对响应因子。该单谱式要求相关直接馈入可测或被独立约束；未观测馈入会带来模型依赖。

对直接馈入目标态的跃迁 $B$，在其飞行移位分量 $B_s$ 上设符合门，设目标态出射跃迁为 $A$。若该门与级联选择纯净，且 $\{B_s,A_s\}$、$\{B_s,A_u\}$ 在各距离采用同一、与距离无关的曝光归一及一致的响应修正，则直接门的 DDCM 关系为~\cite{dewald-2012-ddcm-recoil-distance}
\begin{equation}\label{eq:5-8-ddcm-gated}
\tau_i(x)=\frac{\{B_s,A_u\}(x)}{v\,\mathrm d\{B_s,A_s\}(x)/\mathrm dx}.
\end{equation}
$B_s$ 限定馈入发生在碰到阻止箔之前，馈入可以发生在飞行途中的不同时刻。门选总产额 $\{B_s,A_s\}+\{B_s,A_u\}$ 一般随距离变化，故不能直接以 $-\mathrm d\{B_s,A_u\}/\mathrm dx$ 代替分母中的导数；若逐距离除以门峰面积，求导还须保留归一因子的变化。门污染、未分辨旁路及漏掉的级联需另行评估。对更高处的间接门 $C_s$，若 $B$ 是目标态的直接馈入跃迁，且满足所选级联和响应归一条件，则需扣除馈入项：
\[
\tau_i(x)=\frac{\{C_s,A_u\}(x)-\alpha\{C_s,B_u\}(x)}{v\,\mathrm d\{C_s,A_s\}(x)/\mathrm dx},\qquad
\alpha=\left\langle\frac{\{C_s,A_s\}+\{C_s,A_u\}}{\{C_s,B_s\}+\{C_s,B_u\}}\right\rangle_x.
\]
该式采用距离平均的相对归一因子，其适用性应由数据检验。对离散距离点可采用平滑拟合提取导数，并传播统计及拟合不确定度。在理想恒速模型中，公共零点平移不改变对应点的导数比；真实分析仍须校准相对位移、反冲速度、几何效率及有限速度分布，并评估箔内慢化。

% 第75--87行保留现有力学方程时，应加上非相对论、平均轨迹及平均寿命定义；可在式5-8-f-tau后补充：
式中 $t$ 与速度轨迹均为实验室系量，上式使用低速近似，因而以 $e^{-t/\tau}/\tau$ 近似发射时间概率密度；严格处理时应以本征时间 $s(t)=\int_0^t\mathrm dt'/\gamma(t')$ 替换指数并乘以 $1/\gamma(t)$。对速度及角度分布，应在相应探测效率、角分布和响应加权下计算速度沿探测方向的投影，而非总能将其拆成两个平均因子的乘积。$\theta_{\rm det}=90^\circ$ 时该一阶质心归一式无定义；其他角度或全谱信息能否约束寿命取决于实际响应和统计灵敏度。

% 第92行中“10%--20%”及“直接线性同比例”需要针对具体体系给出证据，建议替换为：
阻止本领模型的不确定度可通过改变模型参数并重新拟合谱形来传播；其对寿命的影响取决于速度历史、馈入及实验响应。

% 图脚注和脚本修复同步：
% 单态指数图无需先归一再求导。DSAM 脚本应在一致能量间隔上卷积，且在整个绘图能域保留展宽，最后仅作标注明确的峰高归一；其余关于峰型和敏感窗的陈述应限定为示意模型，不给普适边界。

% 可在上述时间约定之后加入一个一般前向谱形表达式：
若用 $z$ 汇总初始速度、轨迹与探测方向，$p(z)$ 为其概率密度，$w(t,z)$ 为实验室时间中的发射概率密度，$\epsilon(t,z)$ 为探测权重，且 $R(E\mid E_\gamma,z)$ 在全能量域归一，则预测谱形可写为
\[
Y(E)=C\!\int\!\mathrm dz\,p(z)\!\int_0^\infty\!\mathrm dt\,
 w(t,z)\,\epsilon(t,z)\,R\!\left(E\mid E_\gamma(t,z),z\right).
\]
直接瞬时布居时 $w(t,z)=e^{-s(t,z)/\tau}/[\gamma(t,z)\tau]$；有延迟馈入时应由级联布居方程或馈入时间分布求得 $w$。若仅拟合有限能窗，窗内积分给出该窗接受率；不能先将每个能量分量在窗内独立归一再相加。
